\documentclass[11pt]{article}
\usepackage[margin=1in]{geometry}
\usepackage[T1]{fontenc}
\usepackage{lmodern}
\usepackage{microtype}
\usepackage{xcolor}
\usepackage{hyperref}
\usepackage{booktabs}
\usepackage{enumitem}
\usepackage{listings}
\usepackage{tcolorbox}
\usepackage{parskip}
\usepackage{array}
\newcolumntype{L}[1]{>{\raggedright\arraybackslash}p{#1}}

% \path{} (from url, loaded by hyperref) prints paths and identifiers
% verbatim -- so `--` is not turned into a dash -- and may break at / _ - .
\makeatletter
\g@addto@macro\UrlBreaks{\do\_\do\-\do\.}
\makeatother
\urlstyle{tt}
\setlength{\emergencystretch}{3em}

\hypersetup{
  colorlinks=true,
  linkcolor=blue!60!black,
  urlcolor=blue!60!black
}

\definecolor{codebg}{RGB}{245,245,245}
\definecolor{codecomment}{RGB}{90,100,90}
\definecolor{codekeyword}{RGB}{0,70,140}
\definecolor{codestring}{RGB}{150,60,20}

\lstdefinestyle{pythonstyle}{
  language=Python,
  backgroundcolor=\color{codebg},
  basicstyle=\ttfamily\small,
  keywordstyle=\color{codekeyword}\bfseries,
  commentstyle=\color{codecomment}\itshape,
  stringstyle=\color{codestring},
  showstringspaces=false,
  breaklines=true,
  frame=single,
  rulecolor=\color{black!20},
  columns=fullflexible,
  keepspaces=true
}
\lstdefinestyle{shellstyle}{
  language=bash,
  backgroundcolor=\color{codebg},
  basicstyle=\ttfamily\small,
  showstringspaces=false,
  breaklines=true,
  frame=single,
  rulecolor=\color{black!20},
  columns=fullflexible,
  keepspaces=true
}
\lstset{style=pythonstyle}

\newtcolorbox{keyidea}{
  colback=blue!5, colframe=blue!50!black,
  title=Key Idea, fonttitle=\bfseries
}
\newtcolorbox{tipbox}{
  colback=green!5, colframe=green!45!black,
  title=Tip, fonttitle=\bfseries
}
\newtcolorbox{warningbox}{
  colback=orange!7, colframe=orange!60!black,
  title=Common Pitfall, fonttitle=\bfseries
}

\title{
  \textbf{Codebase Change Summary --- Week 10}\\[4pt]
  \large Integration Testing\\[2pt]
  \normalsize SE 413/513 --- Software Testing
}
\author{Demo Testing Suite}
\date{}

\begin{document}
\maketitle

\begin{center}
\textit{What changed in \path{demo-testing-suite/} since Weeks 8--9, and why}
\end{center}

\vspace{0.3cm}

% ============================================================
\section{Purpose of This Document}
% ============================================================
The \path{demo-testing-suite/} codebase is a single, shared project that
grows one lecture at a time rather than being rebuilt from scratch each
week. This document summarizes exactly what was added or changed for
Week 10 and, more importantly, explains \emph{why} each change exists ---
which concept from the Week 10 reading, \textit{Integration Testing:
Testing How Components Work Together}, it demonstrates. Read this
alongside the reading (\path{reading_materials/w10_integration_testing.pdf})
and \path{demo-testing-suite/README.md}. The reading prints the key code
listings and summarizes the rest; every file it mentions is in
\path{demos/integration_testing/}.

\begin{keyidea}
Nothing from Weeks 1--9 was removed. This week does two things. It adds
a self-contained project that runs every example in the reading: a
small shop with an HTTP API, a database, and a (fake) third-party
payment provider, its interface specifications, a written design
review, and two deliberately broken components. And it applies
integration testing to the \emph{real} suite for the first time ---
which found two real problems there.
\end{keyidea}

% ============================================================
\section{What Changed At a Glance}
% ============================================================
\begin{center}
\small
\begin{tabular}{@{}L{6.0cm}L{1.5cm}L{6.9cm}@{}}
\toprule
\textbf{File} & \textbf{Status} & \textbf{Purpose} \\
\midrule
\multicolumn{3}{@{}l}{\textit{The main suite}}\\
\path{tests/integration/conftest.py} & New & Assembles
\path{UserRegistrationService} from real parts: the real
\path{UserRepository}, a SQLite file, and a fake email client. \\
\path{tests/integration/test_registration.py} & New & Three passing
integration tests, and one strict \texttt{xfail} documenting a defect
they found. \\
\path{app/accounts/user_repository.py} & Modified & Adds
\path{create_schema()}: the production \texttt{users} schema, which
until now existed only in a test fixture. \\
\texttt{pytest.ini} & Modified & Registers the \texttt{integration}
marker. \\
\midrule
\multicolumn{3}{@{}l}{\textit{The reading's examples (not in the default run)}}\\
\path{demos/integration_testing/shop/} & New & The shop: Flask API,
order service, repository, SQLite schema, payment client. \\
\path{demos/integration_testing/shop/payments_v1.py},
\path{service_buggy.py} & New & Two broken components: a client that
sends dollars instead of cents, and a service with a transaction bug. \\
\path{demos/integration_testing/fakes/} & New & A fake payment provider
that enforces the provider's contract, served over real HTTP. \\
\path{demos/integration_testing/specs/} & New & Both API
specifications, and the design/API review (findings R1--R11, each traced
to its tests). \\
\path{demos/integration_testing/tests/} & New & 46 tests: unit,
bottom-up, top-down, stubbed, live, and fully assembled. \\
\path{demos/integration_testing/serve.py}, \texttt{README.md} & New &
Run the shop for \texttt{curl}; every command with its expected output. \\
\midrule
\multicolumn{3}{@{}l}{\textit{Tooling and documentation}}\\
\texttt{requirements.txt} & Modified & Adds \texttt{flask},
\texttt{requests}, and \texttt{responses}. \\
\path{docs/DEMO_GUIDE_week10_integration_testing.md} & New &
Section-by-section lecture companion for the reading. \\
\path{docs/DEMO_GUIDE_week04_tdd_test_doubles.md} & Modified & A
forward note: the Week 4 registration service is now also
integration-tested. \\
\texttt{README.md} (suite and repository) & Modified & Layout,
commands, expected results, and the Weekly Updates index. \\
\bottomrule
\end{tabular}
\end{center}

\begin{tipbox}
Install the new dependencies before running anything this week:
\begin{lstlisting}[style=shellstyle]
cd demo-testing-suite
pip install -r requirements.txt
\end{lstlisting}
\end{tipbox}

% ============================================================
\section{Integration Testing Comes to the Main Suite}
% ============================================================
\textbf{Locations:} \path{tests/integration/},
\path{app/accounts/user_repository.py}

Since Week 4, \path{UserRegistrationService} has been tested with a
\emph{mocked} \path{UserRepository}
(\path{tests/accounts/test_registration_service.py}). Those tests are
still there and still useful. The new folder tests the same service
assembled from real parts, and every test in it is marked:
\begin{lstlisting}
pytestmark = pytest.mark.integration
\end{lstlisting}
\begin{lstlisting}[style=shellstyle]
pytest -m integration          # 3 passed, 75 deselected, 1 xfailed
pytest -m "not integration"    # 1 failed, 72 passed, 4 deselected, 2 xfailed
\end{lstlisting}
The second line is exactly the Weeks 8--9 result.

Writing these tests found two problems, both real:
\begin{enumerate}
\item \textbf{The schema lived only in a test.} The \texttt{users} table
had been created by a test fixture since Week 2, and nowhere in
\path{app/}. An integration test has to use the application's own
schema, so \path{create_schema()} now exists. It also makes
\texttt{name} \texttt{NOT NULL UNIQUE}, so the database rejects a
duplicate even from code that bypasses the service.
\item \textbf{A half-registered user.} \path{register()} saves the user,
\emph{then} sends the welcome email. If the mail server is down, the user
is saved but never welcomed --- and cannot try again, because a second
registration is rejected as a duplicate. A mocked repository has no
saved state to leave behind, so no unit test could see it.
\end{enumerate}

\begin{warningbox}
The second problem is documented with a strict \texttt{xfail}, not
fixed --- the same way Week 1's late-fee bug and Week 3's shipping bug
were. Fixing it is Exercise 4 in the reading. When you do,
\texttt{strict=True} makes the test \emph{fail} (\texttt{XPASS(strict)})
until you remove the marker.
\end{warningbox}

% ============================================================
\section{The Reading's Shop, Runnable}
% ============================================================
\textbf{Location:} \path{demos/integration_testing/}. \textbf{Run its
commands from inside the folder:}
\begin{lstlisting}[style=shellstyle]
cd demos/integration_testing
pytest                         # 46 passed, in about a second
\end{lstlisting}

The shop has six integration points (IP1--IP6): client to API, API to
service, service to repository, repository to SQLite, service to payment
client, and payment client to the third-party provider over HTTP. Each
test file makes a different set of them real:

\begin{center}
\small
\begin{tabular}{@{}L{5.4cm}L{5.0cm}L{3.8cm}@{}}
\toprule
\textbf{Test file} & \textbf{What is real} & \textbf{Strategy step}\\
\midrule
\path{test_order_service_unit.py} & only the service (all mocked) & unit baseline\\
\path{test_repository.py} & repository + SQLite & bottom-up (driver)\\
\path{test_payments_stubbed.py} & client; provider stubbed (\texttt{responses}) & stub at HTTP\\
\path{test_payments_live.py} & client + fake provider, real HTTP & bottom-up\\
\path{test_order_service.py} & service + DB + fake provider & bottom-up, step 2\\
\path{test_api_topdown.py} & API; service stubbed & top-down\\
\path{test_api.py} & everything & assembled system\\
\bottomrule
\end{tabular}
\end{center}

\subsection*{Two broken components, one switch each}
Environment variables swap in the deliberately broken versions, for one
command only. The \emph{v1 client} sends \texttt{"12.50"} dollars instead
of \texttt{1250} cents; the \emph{buggy service} commits the stock and the
order in two separate transactions:
\begin{lstlisting}[style=shellstyle]
SHOP_PAYMENTS=v1   pytest tests/test_api.py
SHOP_SERVICE=buggy pytest tests/test_order_service.py
\end{lstlisting}
\begin{center}
\small
\begin{tabular}{@{}L{4.9cm}L{1.8cm}L{2.3cm}L{5.4cm}@{}}
\toprule
\textbf{Run} & \textbf{v1 client} & \textbf{buggy service} & \textbf{What it shows}\\
\midrule
\path{test_order_service_unit.py} & 3 passed & 3 passed & mocks cannot see either defect\\
\path{test_payments_stubbed.py} & 1 failed & 5 passed & the loose stub passes; the contract check fails\\
\path{test_payments_live.py} & 3 failed & 6 passed & the provider answers \texttt{400}: ``integer number of cents''\\
\path{test_order_service.py} & 2 failed & 1 failed & v1: charge rejected; buggy: \texttt{assert 3 == 5}, stock leaked\\
\path{test_api.py} & 2 failed & 5 passed & \texttt{assert 502 == 201}: \emph{where}?\\
\bottomrule
\end{tabular}
\end{center}

Read the table by column. The dollars-for-cents defect is caught at
every level from the contract check upward, but only the lower-level
tests say \emph{where} it is. The transaction defect is caught at
exactly one level --- the service against a real database --- and the
end-to-end API tests miss it, because the API's input validation blocks
the input that triggers it.

\subsection*{The design review}
\path{specs/design_review.md} records a review held before the tests
were written: the integration points, the checklist, and eleven findings
(R1--R11). Each finding names the tests that check it, and each of those
tests names its finding in its docstring. That traceability is the
pattern the reading asks you to follow in your semester project.

\subsection*{Trying the API by hand}
\begin{lstlisting}[style=shellstyle]
python serve.py                    # shop on :8000, fake provider on :8001
curl -i -X POST http://127.0.0.1:8000/orders -H 'X-API-Key: shop-key' \
     -H 'Content-Type: application/json' \
     -d '{"sku": "MUG", "quantity": 2, "customer_email": "ana@example.com"}'
\end{lstlisting}
returns \texttt{201 CREATED} with a \texttt{Location} header. The folder's
\texttt{README.md} lists the 400, 401, 409, 402, and 503 variations.

% ============================================================
\section{Updated: Expected Test Results}
% ============================================================
From the project root:
\begin{lstlisting}[style=shellstyle]
cd demo-testing-suite
pip install -r requirements.txt
pytest -v
\end{lstlisting}
you should now see \textbf{75 passed, 3 xfailed, 1 failed}: three more
passing tests than Weeks 8--9, the same Week 1 and Week 3 non-passes as
before, and the one new strict \texttt{xfail} described above. Every
\path{app/} module is still at 100\% branch coverage.
\path{demos/integration_testing/} is not part of that run
(\texttt{pytest.ini} sets \texttt{testpaths = tests}); run it from inside
its folder.

% ============================================================
\section{How This Connects to the Reading}
% ============================================================
\begin{keyidea}
Weeks 2--9 tested components one at a time. This week tests the
promises between them: the formats, transactions, credentials, and
failure behavior that two components must agree on. Both demo defects
pass every unit test, and both real defects in the main suite were
invisible to the Week 4 mocks. Each was found by connecting the real
parts, or --- before that --- by reviewing the design.
\end{keyidea}

% ============================================================
\section{Suggested Next Steps}
% ============================================================
\begin{enumerate}
\item Run the switch table above yourself, one cell at a time, and
\emph{predict} each result before you run it.
\item Before reading \path{specs/design_review.md}, review
\path{specs/payment_provider_api.md} yourself with the reading's
checklist. How many of R1--R11 did you find?
\item Fix the half-registered user (Exercise 4). Decide between deleting
the user when the email fails and saving it with a ``welcome pending''
flag. Remove the \texttt{xfail} marker when your fix lands.
\item Implement refunds in the demo (Exercise 3), writing the five
IT-R test cases first.
\item Find out what \path{PaymentClient.charge()} does when the provider
returns \texttt{201} without an \texttt{id} (Exercise 5), and decide
what it should do.
\end{enumerate}

\end{document}
